\section{引言}

图神经网络（GNN）已成为分子图和蛋白质图学习的标准方法，因为它们能够同时编码局部邻域结构与更高阶的关系上下文 \cite{kipf2016semi,gilmer2017neural,hamilton2017inductive,wu2020comprehensive,zhou2020graph}。在生物分子图分类中，这种组合尤为重要，因为判别信号往往同时依赖局部生化基序和更广泛的图级组织。本文聚焦于以 PROTEINS 和 DD 为核心的蛋白质导向 benchmark 设定，将 ENZYMES 作为支持性的生物压力测试，并将 MUTAG、AIDS 和 Mutagenicity 视为补充鲁棒性数据集，而非与前者同等核心的生物证据。因此，本文的目标不是为下游生物发现提供直接验证，而是在蛋白质相关图分类设定下研究残差路由行为。

这一设定下的主要方法论困难在于，更深的消息传递并不会自动带来收益。重复传播可能因过平滑及相关深度效应抑制有信息量的中间状态 \cite{li2018deeper,oono2020simple,chen2023revisiting}，而过压缩则可能将远距离节点信息挤压进无效瓶颈 \cite{alon2020bottleneck,topping2021understanding}。残差连接、归一化与正则化虽能缓解部分退化 \cite{he2016deep,bresson2017residual,zhao2020pair,cai2021graphnorm,rong2019dropedge,chenSimpleDeepGraph2020}，但大多数残差设计仍沿同一路径复用信息，因此它们改善了优化过程，却未直接检验替代性信息流拓扑是否能保留更有用的历史信号。

相关图学习工作已经探索了跳跃连接、稠密聚合、APPNP 式传播、Jumping Knowledge 读出以及层次化池化 \cite{xu2018representation,klicpera2018combining,gao2019graph,ying2018hierarchical}。这些方法能够保留来自不同深度或不同池化层级的信息，但并未直接隔离如下问题：跨分支交换隐藏状态，或跨块边界重新注入池化后的图摘要，是否优于普通的同路径残差复用。此前尚无工作在受控协议下直接比较这些复用拓扑，这正是本文要解决的问题。

本文研究交叉残差图神经网络（CR-GNN），这是一个受控模型家族，在共享图分类主干下比较普通堆叠、节点级残差复用、节点级交叉残差交换、图级残差传播和图级交叉摘要交换。面向论文主叙事的证据主要来自蛋白质导向核心，尤其是 PROTEINS 和 DD；非蛋白质数据集仅用于检验这些结构偏好在核心之外是否仍然稳定。研究围绕三个问题展开：

\begin{enumerate}
    \item \textbf{为什么表示连续性最弱的架构反而最常获胜？} 如果残差复用通常被理解为跨深度保留有效信息的手段，那么当表现最好的架构反而展现出最大的层间表示变化时，这意味着什么？
    \item \textbf{不同复用拓扑何时占优？} 在以 PROTEINS 和 DD 为核心、辅以支持性压力测试和补充鲁棒性数据集的比较中，何时节点级残差复用更强，何时图级残差传播更优，何时交叉残差交换会成为有竞争力甚至更合适的选择？
    \item \textbf{残差路由设计会如何改变答案？} 超越“是否存在额外复用路径”这一二元问题，不同过滤策略，例如可学习的稠密路由、top-k 选择和稀疏抑制，将如何影响残差架构与交叉残差架构的相对表现？
\end{enumerate}

围绕这三个问题组织本文，对界定方法学范围非常关键。本文并不主张某一条额外通路总能带来更好的 GNN，而是试图揭示残差路由在何时、以何种方式有效；普通残差复用何时仍然更强；以及在主干、数据划分和训练协议固定后，为什么层间连续性并不会单调映射到最终准确率。

本文可视为一项面向蛋白质导向生物分子图分类的方法学研究。主要证据来自 PROTEINS 和 DD，ENZYMES 被保留为更困难的支持性生物 benchmark，其较低的绝对准确率本身也有助于刻画任务难度。MUTAG、AIDS 和 Mutagenicity 用于测试蛋白质核心之外的鲁棒性，但它们并不定义本文的主要 biological claim。这种设计使论文保持在具有生物意义的图设定附近，同时仍具备可复现且受控的 benchmark 范围。

\subsection{主要贡献}

本文的主要贡献如下：

\begin{itemize}
    \item \textbf{提出一个用于蛋白质导向图分类复用拓扑比较的受控框架。} 在相同主干、相同数据划分和相同训练协议下，比较五种信息流设计：普通堆叠、节点级残差复用、节点级交叉残差交换、图级残差传播和图级交叉摘要交换。通过固定其他因素，该框架隔离了复用拓扑对性能、一致性和表示动态的影响，其中最强的 biological emphasis 明确落在 PROTEINS 和 DD 上。

    \item \textbf{识别出峰值准确率、排序稳定性和表示连续性之间的系统性权衡。} 没有单一拓扑能同时支配这三个维度。在当前 benchmark 中，图级残差传播最常获得最高峰值准确率，但排序波动明显；交叉残差变体具有竞争力且通常较稳定，但并非普适赢家；而最强的表示连续性并不能预测最佳分类准确率。借助 winner count 和 rank stability 指标量化这种三方张力，为 GNN 架构评价提供了超越单一榜单数字的新视角。

    \item \textbf{给出揭示 CKA--性能悖论的深度扩展机制分析。} 在不同深度上跟踪逐层 CKA、余弦相似度和梯度范数，结果表明表示连续性最弱的架构反而取得了最多获胜次数，这挑战了“更强的层间稳定性意味着更好性能”的假设。这说明当表示变化以结构化方式发生并伴随适当的上下文重新注入时，它可能是有益而非破坏性的。

    \item \textbf{证明路由策略是一个有意义的次级设计轴。} 轻度稀疏路由在代表性目标上最常获胜，尤其常与交叉导向设定配对；而可学习的稠密路由仍在最强的图级残差目标上保持优势。这使讨论从二元架构选择扩展为“拓扑 + 路由”的联合设计问题。
\end{itemize}
